\section{Dense time-varying recursive inner convolution (log-memory baseline)}
\label{sec:inner-dense}

\paragraph{Role in the modular serial concatenation framework.}
In Section~\ref{sec:framework} (the modular serial concatenation framework),
we reduce minimum-distance analysis of the concatenated code
\[
\Csc \ :=\ \Ctwo \circ \Pi \circ \Cone \ \subseteq\ \{0,1\}^N
\]
to bounding, for each outer weight $w$, the inner tail probability
\[
p_w(\delta)\ :=\ 
\PP\!\left[\ \wt\!\big(\Enc_{\mathrm{in}}(U)\big)\le \delta N \ \middle|\ \wt(U)=w\ \right],
\]
where $U$ is uniform over $\{0,1\}^N$ conditioned on $\wt(U)=w$.
(Equivalently: $\Pi$ is uniform, so $\Pi(\text{any fixed weight-}w\text{ word})$ is uniform on the weight-$w$ slice.)

This section instantiates the inner code $\Ctwo$ as a \emph{time-varying dense-tap recursive convolutional encoder}
with memory $m=\Theta(\log N)$, and proves two modular bounds on $p_w(\delta)$:
\begin{enumerate}
\item a \textbf{general three-term tail bound} (useful for large $w$ / large deviations), and
\item a \textbf{forced-termination contraction bound} that converts the ``OFF mechanism'' into a \emph{geometric-in-$w$}
term for small/intermediate $w$.
\end{enumerate}
These are the only inner facts later sections need.

% ============================================================
\subsection{Inner ensemble definition (dense taps, recursive, rate 1)}
\label{subsec:inner-dense-def}

Fix blocklength $N$ and a memory parameter
\[
m \ :=\ \gamma \log_2 N
\qquad\text{for a constant }\gamma>0.
\]
The inner encoder is a streaming linear map $\Enc_{\mathrm{in}}:\F_2^N\to\F_2^N$ defined via a shift-register state.

\paragraph{Streaming definition.}
Initialize $s_1:=0^m\in\F_2^m$.
For each time $t\in[N]$:
\begin{enumerate}
\item Sample a fresh dense tap vector $k_t \gets \Unif(\F_2^m)$, independently over $t$.
\item Output
\begin{equation}
\label{eq:dense-inner-def}
y_t \ :=\ u_t \oplus \langle k_t, s_t\rangle,
\qquad
s_{t+1}\ :=\ \Shift(s_t,y_t),
\end{equation}
where $\Shift$ drops the oldest state bit and appends $y_t$.
\end{enumerate}
Equivalently, for $t\ge 2$, the state stores the last $m$ outputs
\[
s_t \ =\ (y_{t-m},y_{t-m+1},\dots,y_{t-1})
\quad\text{with zero padding for indices }\le 0.
\]

\paragraph{Linearity and generator matrix.}
For a fixed draw of $(k_t)_{t=1}^N$, the map $u\mapsto y$ is $\F_2$-linear and systematic ($y_t$ contains $u_t$),
hence it corresponds to multiplication by an invertible $N\times N$ binary matrix $\Gtwo$:
\[
y \ =\ u \cdot \Gtwo.
\]
We view $\Ctwo$ as the (rate-1) code induced by this random $\Gtwo$ ensemble.

\paragraph{ON/OFF notation.}
Define events
\[
\ON_t \ :=\ \{s_t\neq 0\},
\qquad
\OFF_t\ :=\ \{s_t=0\}.
\]
Intuitively: in \OFF, the encoder outputs $y_t=u_t$ deterministically; in \ON, the term $\langle k_t,s_t\rangle$
is a fresh unbiased bit, so outputs behave like fair coins (up to the fixed $u_t$ mask).

% ============================================================
\subsection{Core ON/OFF lemmas}
\label{subsec:inner-dense-onoff}

Let $\mathcal{F}_{t-1}:=\sigma(u_1^t,k_1^{t-1},y_1^{t-1})$ be the natural past up to time $t-1$.

\begin{lemma}[Fresh fair coin while ON]
\label{lem:dense-fair-while-on}
For every $t$, conditioned on $\mathcal{F}_{t-1}$ and on $\ON_t$, the output bit $y_t$ is uniform:
\[
\PP[y_t=0\mid \mathcal{F}_{t-1},\ON_t]=\PP[y_t=1\mid \mathcal{F}_{t-1},\ON_t]=\tfrac12.
\]
Moreover, for any interval $I=\{a,a+1,\dots,b\}$, conditioned on $\bigcap_{t\in I}\ON_t$,
the bits $\{y_t\}_{t\in I}$ are i.i.d.\ $\mathrm{Bernoulli}(1/2)$.
\end{lemma}

\begin{proof}
Fix $t$. On $\ON_t$, the state $s_t\neq 0$ is fixed given $\mathcal{F}_{t-1}$.
Since $k_t\gets\Unif(\F_2^m)$ and is independent of $\mathcal{F}_{t-1}$,
the dot product $\langle k_t,s_t\rangle$ is a nontrivial linear functional of a uniform vector, hence uniform in $\F_2$.
Therefore $y_t=u_t\oplus \langle k_t,s_t\rangle$ is uniform.

For the joint statement, apply the chain rule over $t=a,\dots,b$; conditioning on the past fixes $s_t$,
and the fresh $k_t$ make each conditional bit uniform, so the joint law is $2^{-(b-a+1)}$.
\end{proof}

\begin{lemma}[ON$\to$OFF requires an $m$-run of zeros in the output]
\label{lem:dense-off-requires-mzeros}
If $\OFF_q$ holds at some time $q\ge 2$, then the preceding $m$ outputs are all zero:
\[
y_{q-m}=y_{q-m+1}=\cdots=y_{q-1}=0.
\]
Equivalently, every ON$\to$OFF transition is witnessed by a contiguous length-$m$ all-zero window in the output stream.
\end{lemma}

\begin{proof}
By construction, $s_q$ is the vector of the last $m$ outputs before time $q$. Thus $s_q=0$ iff those $m$ outputs are all zero.
\end{proof}

% ============================================================
\subsection{A general three-term tail bound for $p_w(\delta)$}
\label{subsec:inner-dense-three-term}

We first prove a bound that is convenient for large $w$ and for the linear-weight regime:
it isolates three distinct mechanisms by which $\wt(Y)$ can be small.

\begin{lemma}[First 1 occurs early for uniform weight-$w$ inputs]
\label{lem:first-one-early}
Let $U$ be uniform over $\{0,1\}^N$ conditioned on $\wt(U)=w\ge 1$.
Let $T:=\min\{t:U_t=1\}$ be the first index containing a $1$.
Then for any $\varepsilon\in(0,1)$,
\[
\PP[T>\varepsilon N]\ \le\ (1-\varepsilon)^w.
\]
\end{lemma}

\begin{proof}
$T>\varepsilon N$ means the first $\lfloor \varepsilon N\rfloor$ coordinates contain no $1$.
Under the uniform weight-$w$ model,
\[
\PP[T>\varepsilon N]
=
\frac{\binom{N-\lfloor\varepsilon N\rfloor}{w}}{\binom{N}{w}}
=
\prod_{i=0}^{w-1}\frac{N-\lfloor\varepsilon N\rfloor-i}{N-i}
\le
\left(\frac{N-\lfloor\varepsilon N\rfloor}{N}\right)^w
\le
(1-\varepsilon)^w.
\]
\end{proof}

\begin{corollary}[Gap survival while ON]
\label{cor:dense-gap-survival}
Fix an interval of $g\ge 0$ consecutive zero inputs, and condition on the event that the encoder is ON at the start of
that interval. Then the probability that the encoder remains ON throughout the whole gap is exactly the probability that
a length-$g$ fair-coin sequence contains no run of $m$ consecutive zeros.
In particular,
\begin{equation}
\label{eq:dense-gap-survival}
\PP[\text{survive a zero gap of length }g \mid \ON]
\;\ge\;
1-g\,2^{-m},
\end{equation}
and if $g<m$ then the survival probability is $1$.
\end{corollary}

\begin{proof}
While the encoder stays ON, Lemma~\ref{lem:dense-fair-while-on} gives i.i.d.\ fair-coin outputs across the gap.
By Lemma~\ref{lem:dense-off-requires-mzeros}, an ON$\to$OFF transition can occur during the gap if and only if one of
those outputs starts a length-$m$ all-zero window. Thus the survival event is exactly the event that the $g$ fair-coin
outputs contain no run of $m$ consecutive zeros. The lower bound \eqref{eq:dense-gap-survival} is the union bound over
the at most $g$ candidate start positions for such a zero window, and for $g<m$ no such window can occur.
\end{proof}

\begin{lemma}[Dense inner: three-term low-weight tail]
\label{lem:dense-three-term-tail}
Let $U$ be uniform over $\{0,1\}^N$ conditioned on $\wt(U)=w\ge 1$, and let $Y=\Enc_{\mathrm{in}}(U)$ be produced by
\eqref{eq:dense-inner-def}. Fix $\delta\in(0,1/2)$ and $\varepsilon\in(0,1)$. Then
\begin{equation}
\label{eq:dense-three-term}
\PP[\wt(Y)\le \delta N \mid \wt(U)=w]
\ \le\
(1-\varepsilon)^w\ +\ N\cdot 2^{-m}\ +\ q_{\delta,\varepsilon}(N),
\end{equation}
where
\[
q_{\delta,\varepsilon}(N)
:=
\PP\!\left[\ \mathrm{Binomial}(\lfloor(1-\varepsilon)N\rfloor,\tfrac12)\ \le\ \delta N\ \right]
\ \le\
2^{-(1-\varepsilon)N\cdot \big(1-H_2(\delta/(1-\varepsilon))\big)}
\]
whenever $\delta\le (1-\varepsilon)/2$.
\end{lemma}

\begin{proof}
Let $T$ be the first index with $U_T=1$. By Lemma~\ref{lem:first-one-early},
\[
\PP[T>\varepsilon N]\le (1-\varepsilon)^w.
\]
On the event $\{T\le \varepsilon N\}$ we have $s_T=0$ (since $s_1=0$ and no earlier $1$ has occurred),
hence $Y_T=U_T=1$, which implies $\ON_{T+1}$.

If the encoder ever transitions from \ON to \OFF after time $T+1$, then by Lemma~\ref{lem:dense-off-requires-mzeros}
there exists a length-$m$ all-zero window in $Y$ starting at some $r\in\{T+1,\dots,N-m+1\}$.
For a fixed $r$, conditioning on $\ON_r$ and using Lemma~\ref{lem:dense-fair-while-on} gives
\[
\PP[Y_r=\cdots=Y_{r+m-1}=0 \mid \ON_r] = 2^{-m},
\]
hence $\PP[Y_r=\cdots=Y_{r+m-1}=0 \ \wedge\ \ON_r]\le 2^{-m}$.
A union bound over all $r$ yields
\[
\PP[\exists\ \text{such window}] \ \le\ N\cdot 2^{-m}.
\]

On the complementary event that $T\le \varepsilon N$ and no ON$\to$OFF transition occurs after $T+1$,
we have $\ON_t$ for all $t\in\{T+1,\dots,N\}$, so Lemma~\ref{lem:dense-fair-while-on} implies
$(Y_{T+1},\dots,Y_N)$ are i.i.d.\ fair coins.
Since $T\le \varepsilon N$ implies $N-T\ge (1-\varepsilon)N$, monotonicity of binomial tails yields the term
$q_{\delta,\varepsilon}(N)$.

Union bounding these three mechanisms proves \eqref{eq:dense-three-term}.
\end{proof}

\paragraph{Interpretation.}
The three terms in \eqref{eq:dense-three-term} correspond to:
(i) the first $1$ appears too late (probability $(1-\varepsilon)^w$),
(ii) an ON$\to$OFF transition occurs (probability $\le N2^{-m}$),
(iii) otherwise, the encoder stays ON for at least $(1-\varepsilon)N$ steps and outputs behave like fair coins, so
$\wt(Y)$ has a large-deviation tail $q_{\delta,\varepsilon}(N)$.

\begin{corollary}[Single-run base factor]
\label{cor:dense-inner-single-run}
Fix $\delta\in(0,1/2)$ and let $U$ be uniform over $\{0,1\}^N$ conditioned on $\wt(U)=1$.
Let $Y=\Enc_{\mathrm{in}}(U)$ be produced by \eqref{eq:dense-inner-def}. Then
\begin{equation}
\label{eq:dense-inner-single-run}
\PP[\wt(Y)\le \delta N \mid \wt(U)=1]
\;\le\;
\frac{\lceil 2\delta N\rceil}{N}
\;+\;
N2^{-m}
\;+\;
\PP[\Binomial(\lceil 2\delta N\rceil,\tfrac12)\le \delta N].
\end{equation}
In particular,
\begin{equation}
\label{eq:dense-inner-single-run-asym}
\PP[\wt(Y)\le \delta N \mid \wt(U)=1]
\;\le\;
2\delta
\;+\;
O\!\left(\frac{1}{N}\right)
\;+\;
N2^{-m}
\;+\;
2^{-\Omega_\delta(N)}.
\end{equation}
\end{corollary}

\begin{proof}
Let $T$ be the unique index with $U_T=1$. Under the uniform weight-$1$ model, $T$ is uniform on $\{1,\dots,N\}$.

On the event that the encoder never transitions from \ON to \OFF after time $T+1$, Lemma~\ref{lem:dense-fair-while-on}
implies that $(Y_{T+1},\dots,Y_N)$ are i.i.d.\ fair coins. If $N-T\ge \lceil 2\delta N\rceil$, then
\[
\PP[\wt(Y)\le \delta N \mid T,\ \text{no OFF after }T+1]
\;\le\;
\PP[\Binomial(\lceil 2\delta N\rceil,\tfrac12)\le \delta N]
\]
by monotonicity in the number of fair-coin steps. Otherwise $N-T<\lceil 2\delta N\rceil$, which is equivalent to
\[
T > N-\lceil 2\delta N\rceil.
\]
Since $T$ is uniform, this exceptional late-start event has probability at most $\lceil 2\delta N\rceil/N$.

Finally, by Lemma~\ref{lem:dense-off-requires-mzeros} and the same union bound as in Lemma~\ref{lem:dense-three-term-tail},
the probability that an ON$\to$OFF transition occurs after time $T+1$ is at most $N2^{-m}$.
Union bounding these three mechanisms proves \eqref{eq:dense-inner-single-run}.

For \eqref{eq:dense-inner-single-run-asym}, note that $\lceil 2\delta N\rceil/N = 2\delta + O(1/N)$, while the
binomial term is exponentially small because $\delta < \lceil 2\delta N\rceil/(2\lceil 2\delta N\rceil)$ for all
sufficiently large $N$.
\end{proof}

% ============================================================
\subsection{Forced-termination contraction (tightening the small-$w$ regime)}
\label{subsec:inner-dense-forced-termination}

For small/intermediate weights, the pointwise polynomial term $N2^{-m}$ in Lemma~\ref{lem:dense-three-term-tail}
is undesirable. We now exploit recursion plus the run structure of a sparse interleaved input to turn the OFF mechanism
into a genuine geometric-in-$w$ contraction.

\begin{definition}[Runs of ones]
\label{def:runs}
For $u\in\{0,1\}^N$, let $R(u)$ denote the number of runs of ones in $u$,
i.e., the number of maximal contiguous substrings of $1$'s.
\end{definition}

\begin{lemma}[Exact run-count distribution under the uniform weight-$w$ model]
\label{lem:run-count-distribution}
Let $U$ be uniform over $\{0,1\}^N$ conditioned on $\wt(U)=w\ge 1$.
Then for every $s\in\{1,\dots,w\}$,
\[
\PP[R(U)=s]
\;=\;
\frac{\binom{w-1}{s-1}\binom{N-w+1}{s}}{\binom{N}{w}}.
\]
Consequently, for every $r\in\{1,\dots,w\}$,
\begin{equation}
\label{eq:run-lower-tail-exact}
\PP[R(U)<r]
\;=\;
\frac{1}{\binom{N}{w}}
\sum_{s=1}^{r-1}
\binom{w-1}{s-1}\binom{N-w+1}{s}.
\end{equation}
\end{lemma}

\begin{proof}
A binary word of length $N$ with weight $w$ and exactly $s$ runs of ones is described by two independent choices.
First choose the $s$ positive run lengths whose sum is $w$, which gives $\binom{w-1}{s-1}$ possibilities.
Next place the $N-w$ zeros into the $s+1$ zero-gaps around those runs, with the $s-1$ internal gaps required to
receive at least one zero so that adjacent runs remain distinct.
After reserving one zero for each internal gap, the remaining $N-w-(s-1)$ zeros are distributed freely over the
$s+1$ gaps, giving
\[
\binom{N-w-(s-1)+s}{s}
\;=\;
\binom{N-w+1}{s}
\]
choices.
Dividing by the total number $\binom{N}{w}$ of weight-$w$ words yields the stated probability formula,
and summing over $s<r$ gives \eqref{eq:run-lower-tail-exact}.
\end{proof}

\begin{lemma}[Gap composition conditioned on weight and run count]
\label{lem:gap-composition-given-runs}
Let $U$ be uniform over $\{0,1\}^N$ conditioned on $\wt(U)=w\ge 1$ and $R(U)=s$.
Write the zero gaps around the $s$ runs as
\[
z_0,\ z_1,\ \dots,\ z_{s-1},\ z_s,
\]
where $z_0$ and $z_s$ are the leading and trailing zero counts and each internal gap $z_i$ for
$1\le i\le s-1$ is at least $1$. Define shifted internal gaps
\[
y_0:=z_0,\qquad y_i:=z_i-1 \ (1\le i\le s-1),\qquad y_s:=z_s.
\]
Then $(y_0,\dots,y_s)$ is uniform over all weak compositions of
\[
M:=N-w-(s-1)
\]
into $s+1$ parts.

In particular, for any fixed internal gap $G=z_i$ with $1\le i\le s-1$,
\begin{equation}
\label{eq:internal-gap-marginal}
\PP[G=g \mid \wt(U)=w,\ R(U)=s]
\;=\;
\frac{\binom{M-g+s-1}{s-1}}{\binom{M+s}{s}}
\qquad (g\ge 1),
\end{equation}
whenever $M-g\ge 0$, and therefore
\begin{equation}
\label{eq:internal-gap-mean}
\EE[G \mid \wt(U)=w,\ R(U)=s]
\;=\;
1+\frac{M}{s+1}
\;=\;
\frac{N-w+2}{s+1}.
\end{equation}
\end{lemma}

\begin{proof}
Conditioned on $\wt(U)=w$ and $R(U)=s$, a word is specified by:
\begin{itemize}
\item a composition of $w$ into $s$ positive run lengths, and
\item a choice of zero gaps $z_0,z_1,\dots,z_s$ with $z_0,z_s\ge 0$, $z_i\ge 1$ for $1\le i\le s-1$, and
\(\sum_{i=0}^s z_i=N-w\).
\end{itemize}
After the shift $y_i=z_i-1$ on the internal gaps, the second datum becomes a weak composition of
$M=N-w-(s-1)$ into $s+1$ parts. For each fixed shifted gap vector there are exactly
$\binom{w-1}{s-1}$ compatible run-length compositions, so under the uniform law on the $(w,s)$-slice the shifted
gap vector is itself uniform over those weak compositions.

Fix an internal gap $G=z_i$ and write $G=1+Y_i$. The event $\{G=g\}$ is equivalent to $Y_i=g-1$, leaving
$M-(g-1)=M-g+1$ units to be distributed across the remaining $s$ coordinates. Hence
\[
\PP[G=g \mid \wt(U)=w,\ R(U)=s]
\;=\;
\frac{\binom{M-g+s-1}{s-1}}{\binom{M+s}{s}},
\]
which is \eqref{eq:internal-gap-marginal}. Finally, by symmetry of the uniform composition law,
$\EE[Y_i]=M/(s+1)$, so \eqref{eq:internal-gap-mean} follows from $G=1+Y_i$.
\end{proof}

\begin{corollary}[Forced-cluster count from short gaps]
\label{cor:forced-cluster-count}
Fix a memory parameter $m\ge 1$. Under the same conditioning as
Lemma~\ref{lem:gap-composition-given-runs}, define the \emph{forced cluster count}
\[
C_m(U)
\;:=\;
1+\#\{\,i\in\{1,\dots,s-1\}: z_i\ge m\,\}.
\]
Equivalently, one merges adjacent runs whenever their separating zero gap is $<m$.
Then
\begin{equation}
\label{eq:forced-cluster-prob}
\PP[z_i\ge m \mid \wt(U)=w,\ R(U)=s]
\;=\;
\frac{\binom{M-m+s+1}{s}}{\binom{M+s}{s}}
\end{equation}
for each internal gap \(z_i\), whenever \(M-m+1\ge 0\), and therefore
\begin{equation}
\label{eq:forced-cluster-mean}
\EE[C_m(U)\mid \wt(U)=w,\ R(U)=s]
\;=\;
1+(s-1)\frac{\binom{M-m+s+1}{s}}{\binom{M+s}{s}}.
\end{equation}
\end{corollary}

\begin{proof}
By definition, \(z_i\ge m\) if and only if \(y_i=z_i-1\ge m-1\).
Under Lemma~\ref{lem:gap-composition-given-runs}, the shifted gap vector is uniform over weak compositions of
\(M\) into \(s+1\) parts. So fixing \(y_i\ge m-1\) and subtracting \(m-1\) from that coordinate leaves a weak
composition of \(M-(m-1)\) into \(s+1\) parts. Hence
\[
\PP[z_i\ge m \mid \wt(U)=w,\ R(U)=s]
\;=\;
\frac{\binom{M-(m-1)+s}{s}}{\binom{M+s}{s}}
\;=\;
\frac{\binom{M-m+s+1}{s}}{\binom{M+s}{s}},
\]
which is \eqref{eq:forced-cluster-prob}. The expectation formula
\eqref{eq:forced-cluster-mean} then follows by linearity of expectation over the \(s-1\) exchangeable internal gaps.
\end{proof}

\begin{lemma}[First-run late tail conditioned on the run count]
\label{lem:first-run-late-given-runs}
Let $U$ be uniform over $\{0,1\}^N$ conditioned on $\wt(U)=w\ge 1$ and $R(U)=s$.
Let $T_1(U)$ denote the start position of the first run of ones.
Then for every integer $t\in\{0,1,\dots,N-w\}$,
\begin{equation}
\label{eq:first-run-late-given-runs}
\PP[T_1(U)>t \mid \wt(U)=w,\ R(U)=s]
\;=\;
\frac{\binom{N-t-w+1}{s}}{\binom{N-w+1}{s}}.
\end{equation}
Consequently,
\begin{equation}
\label{eq:first-run-late-given-runs-simple}
\PP[T_1(U)>t \mid \wt(U)=w,\ R(U)=s]
\;\le\;
\left(1-\frac{t}{N}\right)^s.
\end{equation}
\end{lemma}

\begin{proof}
The event $\{T_1(U)>t\}$ means that the first $t$ coordinates are all zero, so the entire weight-$w$, $s$-run pattern
must be placed in the remaining $N-t$ coordinates. By Lemma~\ref{lem:run-count-distribution}, the number of such words is
\[
\binom{w-1}{s-1}\binom{N-t-w+1}{s},
\]
while the total number of weight-$w$, $s$-run words is
\[
\binom{w-1}{s-1}\binom{N-w+1}{s}.
\]
Taking the ratio gives \eqref{eq:first-run-late-given-runs}.

For the bound, write
\[
\frac{\binom{N-t-w+1}{s}}{\binom{N-w+1}{s}}
\;=\;
\prod_{j=0}^{s-1}\frac{N-t-w+1-j}{N-w+1-j}
\;=\;
\prod_{j=0}^{s-1}\left(1-\frac{t}{N-w+1-j}\right).
\]
Since $N-w+1-j\le N$, each factor is at most $1-t/N$, proving
\eqref{eq:first-run-late-given-runs-simple}.
\end{proof}

\begin{corollary}[Conditioned late-start base factor]
\label{cor:first-run-late-qbase}
Fix $\delta\in(0,1/2)$ and let
\[
L:=\lceil 2\delta N\rceil.
\]
Under the same conditioning as Lemma~\ref{lem:first-run-late-given-runs},
\begin{equation}
\label{eq:first-run-late-qbase}
\PP[T_1(U)>N-L \mid \wt(U)=w,\ R(U)=s]
\;=\;
\frac{\binom{L-w+1}{s}}{\binom{N-w+1}{s}}
\;\le\;
\left(\frac{L}{N}\right)^s.
\end{equation}
In particular,
\begin{equation}
\label{eq:first-run-late-qbase-asym}
\PP[T_1(U)>N-L \mid \wt(U)=w,\ R(U)=s]
\;\le\;
\left(2\delta + O\!\left(\frac{1}{N}\right)\right)^s.
\end{equation}
\end{corollary}

\begin{proof}
Apply Lemma~\ref{lem:first-run-late-given-runs} with $t=N-L$.
Then
\[
\PP[T_1(U)>N-L \mid \wt(U)=w,\ R(U)=s]
\;=\;
\frac{\binom{L-w+1}{s}}{\binom{N-w+1}{s}}.
\]
The bound follows from
\[
\frac{\binom{L-w+1}{s}}{\binom{N-w+1}{s}}
\;=\;
\prod_{j=0}^{s-1}\frac{L-w+1-j}{N-w+1-j}
\;\le\;
\left(\frac{L}{N}\right)^s,
\]
and $L/N = 2\delta + O(1/N)$ gives \eqref{eq:first-run-late-qbase-asym}.
\end{proof}

\begin{corollary}[Two isolated ones: exact late-placement factor]
\label{cor:first-run-late-two-run}
Fix $\delta\in(0,1/2)$ and let $L=\lceil 2\delta N\rceil$.
If $U$ is uniform over $\{0,1\}^N$ conditioned on $\wt(U)=2$ and $R(U)=2$, then
\begin{equation}
\label{eq:first-run-late-two-run}
\PP[T_1(U)>N-L \mid \wt(U)=2,\ R(U)=2]
\;=\;
\frac{\binom{L-1}{2}}{\binom{N-1}{2}}.
\end{equation}
Equivalently,
\begin{equation}
\label{eq:first-run-late-two-run-asym}
\PP[T_1(U)>N-L \mid \wt(U)=2,\ R(U)=2]
\;=\;
\left(\frac{L}{N}\right)^2
\left(1+O\!\left(\frac{1}{L}\right)\right).
\end{equation}
\end{corollary}

\begin{proof}
This is the special case $(w,s)=(2,2)$ of Corollary~\ref{cor:first-run-late-qbase}. The asymptotic form follows from
\[
\frac{\binom{L-1}{2}}{\binom{N-1}{2}}
\;=\;
\left(\frac{L}{N}\right)^2
\frac{(1-1/L)(1-2/L)}{(1-1/N)(1-2/N)}.
\]
\end{proof}

\begin{corollary}[Isolated order statistics: exact late tail]
\label{cor:isolated-order-late}
Fix integers \(1\le r\le h\le N\) and let \(L\in\{1,\dots,N\}\).
Let \(U\) be uniform over \(\{0,1\}^N\) conditioned on \(\wt(U)=h\) and \(R(U)=h\), and write
\[
X_1(U)<X_2(U)<\cdots<X_h(U)
\]
for the ordered one positions.
Then
\begin{equation}
\label{eq:isolated-order-late}
\PP[X_r(U)>N-L \mid \wt(U)=h,\ R(U)=h]
\;=\;
\frac{1}{\binom{N-h+1}{h}}
\sum_{j=0}^{r-1}
\binom{N-L-r+1}{j}\binom{L-h+r}{h-j},
\end{equation}
with the convention that the sum is interpreted as \(1\) when \(N-L-r+1<0\).
\end{corollary}

\begin{proof}
Conditioned on \(\wt(U)=h\) and \(R(U)=h\), the isolated one positions satisfy
\[
1\le X_1<\cdots<X_h\le N,
\qquad
X_{i+1}\ge X_i+2.
\]
Define
\[
Y_i:=X_i-i+1,\qquad i=1,\dots,h.
\]
Then
\[
1\le Y_1<\cdots<Y_h\le N-h+1,
\]
and the map \(X\mapsto Y\) is a bijection from isolated \(h\)-subsets of \([N]\) to ordinary \(h\)-subsets of
\([N-h+1]\). Hence \(Y_1<\cdots<Y_h\) is uniform over all \(h\)-subsets of \([N-h+1]\).

Now \(X_r>N-L\) is equivalent to
\[
Y_r=X_r-r+1 > N-L-r+1.
\]
So among the \(Y_i\), at most \(r-1\) may lie in the prefix \(\{1,\dots,N-L-r+1\}\), while the remaining
\(h-j\) elements must lie in the tail of size
\[
(N-h+1)-(N-L-r+1)=L-h+r.
\]
Choosing exactly \(j\) prefix elements gives
\[
\binom{N-L-r+1}{j}\binom{L-h+r}{h-j},
\]
and summing over \(j=0,\dots,r-1\) proves \eqref{eq:isolated-order-late}.
\end{proof}

\begin{corollary}[Isolated early-count decomposition]
\label{cor:isolated-early-count}
Fix \(h\le N\) and let \(L\in\{1,\dots,N\}\).
Let \(U\) be uniform over \(\{0,1\}^N\) conditioned on \(\wt(U)=h\) and \(R(U)=h\), and define
\[
J_L(U)
\;:=\;
\#\{\,i\in\{1,\dots,h\}: X_i(U)\le N-L\,\},
\]
the number of isolated one-positions that start before the final \(L\) coordinates.
Then for every \(j\in\{0,\dots,h\}\),
\begin{equation}
\label{eq:isolated-early-count}
\PP[J_L(U)=j \mid \wt(U)=h,\ R(U)=h]
\;=\;
\frac{\binom{N-L}{j}\binom{L-h+1}{h-j}}{\binom{N-h+1}{h}},
\end{equation}
with the convention that the right-hand side is \(0\) when either binomial coefficient is undefined.
\end{corollary}

\begin{proof}
Under the shift \(Y_i=X_i-i+1\) from the proof of Corollary~\ref{cor:isolated-order-late},
the isolated slice \(\wt(U)=h,\ R(U)=h\) becomes a uniform \(h\)-subset
\[
1\le Y_1<\cdots<Y_h\le N-h+1.
\]
The condition \(X_i\le N-L\) is equivalent to \(Y_i\le N-L-i+1\), and therefore
\(J_L(U)=j\) means exactly \(j\) of the shifted positions lie in the prefix \(\{1,\dots,N-L\}\),
while the remaining \(h-j\) lie in the tail of size
\[
(N-h+1)-(N-L)=L-h+1.
\]
Choosing those two groups independently gives
\[
\binom{N-L}{j}\binom{L-h+1}{h-j}
\]
valid isolated \(h\)-subsets. Dividing by the total \(\binom{N-h+1}{h}\) proves
\eqref{eq:isolated-early-count}.
\end{proof}

\begin{corollary}[Expected number of early isolated pairs]
\label{cor:isolated-early-pair-mean}
Under the same conditioning and notation as Corollary~\ref{cor:isolated-early-count},
\begin{equation}
\label{eq:isolated-early-pair-mean}
\EE\!\left[\binom{J_L(U)}{2}\,\middle|\, \wt(U)=h,\ R(U)=h\right]
\;=\;
\binom{h}{2}\,
\frac{\binom{N-L}{2}}{\binom{N-h+1}{2}}.
\end{equation}
\end{corollary}

\begin{proof}
Under the shift \(Y_i=X_i-i+1\) from the proof of Corollary~\ref{cor:isolated-order-late}, the isolated slice becomes a
uniform \(h\)-subset of \([N-h+1]\), and \(J_L(U)\) counts how many of those \(h\) points land in the distinguished
prefix of size \(N-L\). So \(J_L(U)\) has the usual hypergeometric distribution with population size \(N-h+1\),
success population \(N-L\), and sample size \(h\). The factorial second moment of a hypergeometric variable gives
\[
\EE[J_L(U)(J_L(U)-1)\mid \wt(U)=h,\ R(U)=h]
\;=\;
h(h-1)\frac{(N-L)(N-L-1)}{(N-h+1)(N-h)}.
\]
Dividing by \(2\) proves \eqref{eq:isolated-early-pair-mean}.
\end{proof}

\begin{corollary}[Tiny weights are almost all isolated]
\label{cor:run-count-full}
Let $U$ be uniform over $\{0,1\}^N$ conditioned on $\wt(U)=h\ge 1$.
Then
\begin{equation}
\label{eq:run-count-full}
\PP[R(U)=h]
\;=\;
\frac{\binom{N-h+1}{h}}{\binom{N}{h}}.
\end{equation}
In particular, for every fixed $h$,
\begin{equation}
\label{eq:run-count-full-asym}
\PP[R(U)=h]
\;=\;
1-O_h\!\left(\frac{1}{N}\right),
\end{equation}
and more concretely,
\begin{equation}
\label{eq:run-count-full-lower}
\PP[R(U)=h]
\;\ge\;
1-\frac{h(h-1)}{N}.
\end{equation}
\end{corollary}

\begin{proof}
The exact formula is the case $s=h$ of Lemma~\ref{lem:run-count-distribution}. For the lower bound, let
\[
A_i:=\{U_i=U_{i+1}=1\},\qquad i=1,\dots,N-1.
\]
If $R(U)<h$, then at least one adjacency event $A_i$ occurs. Under the uniform weight-$h$ model,
\[
\PP[A_i]
\;=\;
\frac{\binom{N-2}{h-2}}{\binom{N}{h}}
\;=\;
\frac{h(h-1)}{N(N-1)}.
\]
Hence by a union bound,
\[
\PP[R(U)<h]
\;\le\;
(N-1)\frac{h(h-1)}{N(N-1)}
\;=\;
\frac{h(h-1)}{N},
\]
which proves \eqref{eq:run-count-full-lower}. The asymptotic claim \eqref{eq:run-count-full-asym} is immediate for
fixed \(h\).
\end{proof}

\begin{lemma}[Late-start plus one-episode correction]
\label{lem:dense-late-start-one-episode}
Let $U$ be uniform over $\{0,1\}^N$ conditioned on $\wt(U)=w\ge 1$ and $R(U)=s$.
Let $Y=\Enc_{\mathrm{in}}(U)$, fix an integer $L$ with $1\le L\le N$, and let $T_1(U)$ be the first run start.
Then for every threshold $d\ge 0$,
\begin{equation}
\label{eq:dense-late-start-one-episode}
\PP[\wt(Y)\le d \mid \wt(U)=w,\ R(U)=s]
\;\le\;
\PP[T_1(U)>N-L \mid \wt(U)=w,\ R(U)=s]
\;+\;
L2^{-m}
\;+\;
\PP[\Binomial(L,\tfrac12)\le d].
\end{equation}
\end{lemma}

\begin{proof}
Let $E_{\mathrm{late}}:=\{T_1(U)>N-L\}$.
On the complementary event $E_{\mathrm{late}}^c$, the first run starts by time $N-L$.
As in the proof of Corollary~\ref{cor:dense-inner-single-run}, that first run forces the encoder into the ON state at
the next time step, thereby initiating the first ON-episode.

If that first ON-episode terminates within its first $L$ steps, then by Lemma~\ref{lem:dense-off-requires-mzeros} and
Lemma~\ref{lem:dense-fair-while-on}, some length-$m$ all-zero output window must begin within those $L$ steps. By a
union bound this has probability at most $L2^{-m}$.

Otherwise the first ON-episode survives for at least $L$ steps. On that whole interval the outputs are i.i.d.\ fair
coins by Lemma~\ref{lem:dense-fair-while-on}. Since the total weight of $Y$ is at most $d$, in particular those first
$L$ output bits of the episode must contain at most $d$ ones. The probability of that event is exactly
\(\PP[\Binomial(L,\tfrac12)\le d]\).

Thus, on $E_{\mathrm{late}}^c$, the low-weight event is contained in the union of the ``early termination within $L$''
event and the length-$L$ binomial tail event. Adding back the late-start event $E_{\mathrm{late}}$ proves
\eqref{eq:dense-late-start-one-episode}.
\end{proof}

\begin{corollary}[Run-mixture slice bound from one episode]
\label{cor:dense-late-start-slice}
Fix $\delta\in(0,1/2)$ and let $d=\lfloor \delta N\rfloor$.
Then for every $w\ge 1$,
\begin{equation}
\label{eq:dense-late-start-slice}
p_w(\delta)
\;\le\;
\sum_{s=1}^{w}
\PP[R(U)=s]\,
\frac{\binom{L-w+1}{s}}{\binom{N-w+1}{s}}
\;+\;
L2^{-m}
\;+\;
\PP[\Binomial(L,\tfrac12)\le d],
\end{equation}
where \(U\) is uniform on the weight-\(w\) slice and \(L\in\{1,\dots,N\}\) is arbitrary.
\end{corollary}

\begin{proof}
Average Lemma~\ref{lem:dense-late-start-one-episode} over the conditional run-count law from
Lemma~\ref{lem:run-count-distribution}, then substitute Corollary~\ref{cor:first-run-late-qbase} for the
late-start term.
\end{proof}

\begin{lemma}[Entropy bound for the run lower tail]
\label{lem:run-lower-tail-entropy}
Fix constants \(0<\omega<1\) and \(0<\rho<\min(\omega,1-\omega)\), and let
\[
w=\lfloor \omega N\rfloor,
\qquad
r=\lfloor \rho N\rfloor.
\]
Assume \(\rho<\omega(1-\omega)\). Then
\begin{equation}
\label{eq:run-lower-tail-entropy}
\PP[R(U)<r]
\;\le\;
N^{O(1)}\,2^{-N\Psi_{\mathrm{run}}(\omega,\rho)},
\end{equation}
where
\begin{equation}
\label{eq:run-rate-function}
\Psi_{\mathrm{run}}(\omega,\rho)
\;:=\;
H_2(\omega)
- \omega H_2\!\Bigl(\frac{\rho}{\omega}\Bigr)
- (1-\omega)H_2\!\Bigl(\frac{\rho}{1-\omega}\Bigr).
\end{equation}
\end{lemma}

\begin{proof}
By \eqref{eq:run-lower-tail-exact},
\[
\PP[R(U)<r]
\;=\;
\frac{1}{\binom{N}{w}}
\sum_{s=1}^{r-1}
\binom{w-1}{s-1}\binom{N-w+1}{s}.
\]
For \(s=\lfloor \sigma N\rfloor\) with \(0<\sigma<\rho\), Stirling's formula gives
\[
\log_2 \binom{w-1}{s-1}
=
\omega N\,H_2\!\Bigl(\frac{\sigma}{\omega}\Bigr) + O(\log N),
\]
\[
\log_2 \binom{N-w+1}{s}
=
(1-\omega)N\,H_2\!\Bigl(\frac{\sigma}{1-\omega}\Bigr) + O(\log N),
\]
and
\[
\log_2 \binom{N}{w}
=
N H_2(\omega) + O(\log N).
\]
Thus each summand has the form
\[
N^{O(1)}\,2^{-N\Psi_{\mathrm{run}}(\omega,\sigma)}.
\]
Moreover,
\[
\frac{\partial}{\partial \sigma}\Psi_{\mathrm{run}}(\omega,\sigma)
=
\log_2\!\left(\frac{(\omega-\sigma)(1-\omega-\sigma)}{\sigma^2}\right),
\]
which is positive throughout the lower-tail regime \(\sigma<\omega(1-\omega)\).
Hence \(\Psi_{\mathrm{run}}(\omega,\sigma)\) is increasing on \((0,\rho)\), so the largest summand occurs at the boundary
\(\sigma\approx \rho\).
Since the sum contains at most \(N\) terms, \eqref{eq:run-lower-tail-entropy} follows.
\end{proof}

\begin{corollary}[Run lower-tail exponent with \(r\) proportional to \(w\)]
\label{cor:run-lower-tail-alpha}
Fix constants \(0<\omega<1\) and \(0<\alpha<1-\omega\), and let
\[
w=\lfloor \omega N\rfloor,
\qquad
r=\lfloor \alpha w\rfloor.
\]
Then
\begin{equation}
\label{eq:run-lower-tail-alpha}
\PP[R(U)<r]
\;\le\;
N^{O(1)}\,2^{-N\Psi_{\mathrm{run}}^{(w)}(\omega,\alpha)},
\end{equation}
where
\begin{equation}
\label{eq:run-rate-function-alpha}
\Psi_{\mathrm{run}}^{(w)}(\omega,\alpha)
\;:=\;
H_2(\omega)
- \omega H_2(\alpha)
- (1-\omega)H_2\!\Bigl(\frac{\alpha\omega}{1-\omega}\Bigr).
\end{equation}
\end{corollary}

\begin{proof}
Apply Lemma~\ref{lem:run-lower-tail-entropy} with \(\rho=\alpha\omega\).
The condition \(\alpha<1-\omega\) is exactly \(\rho<\omega(1-\omega)\).
\end{proof}

\begin{lemma}[Forced-termination contraction]
\label{lem:dense-forced-termination}
Fix $\delta\in(0,1/2)$ and choose any $\xi>0$. Let
\[
L\ :=\ \left\lceil (2+\xi)\delta N \right\rceil.
\]
Let $U$ be uniform over $\{0,1\}^N$ conditioned on $\wt(U)=w\ge 1$, and let $Y=\Enc_{\mathrm{in}}(U)$ be produced by
\eqref{eq:dense-inner-def}. Then for every integer $r\in\{1,\dots,w\}$,
\begin{equation}
\label{eq:dense-forced-termination}
\PP[\wt(Y)\le \delta N \mid \wt(U)=w]
\ \le\
\PP[R(U)<r]\ +\ (L2^{-m})^{r}\ +\ r\cdot \PP[\mathrm{Binomial}(L,\tfrac12)\le \delta N].
\end{equation}
Moreover,
\begin{equation}
\label{eq:runs-adjacent-bound}
\PP[R(U)<w]\ \le\ \frac{w(w-1)}{N-1}.
\end{equation}
In particular, for $w\ll \sqrt{N}$ (e.g.\ $w=O(\log N)$), taking $r=w$ yields a geometric-in-$w$ term
$(L2^{-m})^w$ up to a negligible $O(w^2/N)$ correction.
\end{lemma}

\begin{proof}
\emph{Step 1 (runs initiate ON-episodes).}
When the encoder is OFF ($s_t=0$), we have $y_t=u_t$.
Therefore, at the first position of any run of ones in $u$, the encoder outputs $1$, and the next state becomes nonzero,
so the encoder enters ON immediately after the run begins.
Thus each run of ones initiates an ON-episode.

\emph{Step 2 (small output weight forces fast termination or a binomial deviation).}
Condition on the event $\{R(U)\ge r\}$ and select any $r$ distinct runs.
Consider the corresponding $r$ ON-episodes (they occur on disjoint time intervals since the runs are disjoint).

If any selected ON-episode has length at least $L$, then on the sub-interval where the encoder stays ON,
Lemma~\ref{lem:dense-fair-while-on} gives fair-coin outputs.
Since $L=(2+\xi)\delta N>2\delta N$, the probability that this single episode contributes at most $\delta N$ ones is
upper bounded by $\PP[\mathrm{Binomial}(L,1/2)\le \delta N]$ (monotonicity).
By a union bound over the $r$ selected episodes, the probability that some selected episode is long and still contributes
$\le \delta N$ ones is at most
\[
r\cdot \PP[\mathrm{Binomial}(L,1/2)\le \delta N].
\]

Otherwise, all $r$ selected ON-episodes have length $<L$.
An ON-episode ends only via an ON$\to$OFF transition, which by Lemma~\ref{lem:dense-off-requires-mzeros} requires a
length-$m$ all-zero window in the output while ON.
For a given ON-episode to terminate within $L$ steps, there are at most $L$ candidate start indices for such a window;
for each fixed candidate, Lemma~\ref{lem:dense-fair-while-on} gives probability exactly $2^{-m}$ for an all-zero window.
Thus, by a union bound, a fixed ON-episode terminates within $L$ steps with probability at most $L2^{-m}$.
Since tap vectors are independent over time, the termination events for the $r$ disjoint ON-episodes are independent
conditional on the input pattern. Hence the probability that all $r$ selected ON-episodes terminate within $L$ is
at most $(L2^{-m})^r$.

Finally, union bound over the complementary event $\{R(U)<r\}$ yields \eqref{eq:dense-forced-termination}.

\emph{Step 3 (adjacent-ones bound).}
We have $R(U)<w$ iff $U$ contains an adjacent pair of ones.
Let
\[
A:=\sum_{i=1}^{N-1}\mathbf{1}\{U_i=U_{i+1}=1\}.
\]
Then $\{R(U)<w\}\subseteq\{A\ge 1\}$, so $\PP[R(U)<w]\le \EE[A]$.
For uniform weight-$w$ vectors,
\[
\PP[U_i=U_{i+1}=1]=\frac{\binom{N-2}{w-2}}{\binom{N}{w}}=\frac{w(w-1)}{N(N-1)}.
\]
Therefore $\EE[A]=(N-1)\cdot \frac{w(w-1)}{N(N-1)}=\frac{w(w-1)}{N}$, which implies \eqref{eq:runs-adjacent-bound}
(up to the minor $N$ vs.\ $N-1$ normalization).
\end{proof}

% ============================================================
\subsection{Inner interface theorem (what the outer+interleaver proof consumes)}
\label{subsec:inner-dense-interface}

We package the previous lemmas into an explicit interface for Section~\ref{sec:framework}.

\begin{theorem}[Dense inner interface for the framework]
\label{thm:inner-dense-interface}
Let $\Ctwo$ be the dense time-varying recursive inner encoder \eqref{eq:dense-inner-def} with memory $m$.
Fix $\delta\in(0,1/2)$.
For each $w\ge 1$, let $U$ be uniform over $\{0,1\}^N$ conditioned on $\wt(U)=w$ and let $Y=\Enc_{\mathrm{in}}(U)$.

Then:
\begin{enumerate}
\item (\textbf{General tail, all $w$}) For every $\varepsilon\in(0,1)$,
\[
p_w(\delta)\ :=\ \PP[\wt(Y)\le \delta N\mid \wt(U)=w]
\ \le\
(1-\varepsilon)^w\ +\ N2^{-m}\ +\ q_{\delta,\varepsilon}(N),
\]
with $q_{\delta,\varepsilon}(N)$ as in Lemma~\ref{lem:dense-three-term-tail}.
\item (\textbf{Forced-termination contraction, tunable in $r$}) For every $\xi>0$ and
$L=\lceil(2+\xi)\delta N\rceil$, and for every $r\in\{1,\dots,w\}$,
\[
p_w(\delta)
\ \le\
\PP[R(U)<r]\ +\ (L2^{-m})^{r}\ +\ r\cdot \PP[\mathrm{Binomial}(L,\tfrac12)\le \delta N].
\]
\end{enumerate}
\end{theorem}

\begin{corollary}[Optimization-ready dense-inner envelopes]
\label{cor:dense-inner-piecewise}
Fix $\delta\in(0,1/2)$.
Choose any $\varepsilon\in(0,1)$ and any $\xi>0$, and define
\[
L:=\lceil(2+\xi)\delta N\rceil,
\]
\[
q_{\delta,\varepsilon}(N)
:=
\PP\!\left[\ \Binomial(\lfloor(1-\varepsilon)N\rfloor,\tfrac12)\le \delta N\ \right],
\]
\[
b_{\delta,\xi}(N)
:=
\PP\!\left[\ \Binomial(L,\tfrac12)\le \delta N\ \right].
\]
For $1\le r\le w$, define the exact run lower-tail
\[
\mathrm{RunTail}_N(w,r)
\;:=\;
\frac{1}{\binom{N}{w}}
\sum_{s=1}^{r-1}
\binom{w-1}{s-1}\binom{N-w+1}{s}.
\]
Then for every $w\ge 1$,
\begin{equation}
\label{eq:dense-inner-global-envelope}
p_w(\delta)
\;\le\;
(1-\varepsilon)^w + N2^{-m} + q_{\delta,\varepsilon}(N),
\end{equation}
and for every $r\in\{1,\dots,w\}$,
\begin{equation}
\label{eq:dense-inner-run-envelope}
p_w(\delta)
\;\le\;
\mathrm{RunTail}_N(w,r) + (L2^{-m})^r + r\,b_{\delta,\xi}(N).
\end{equation}
In particular,
\begin{equation}
\label{eq:dense-inner-run-opt-envelope}
p_w(\delta)
\;\le\;
B_w^{\mathrm{run}}(\delta;\xi,m)
\;:=\;
\min_{1\le r\le w}
\left(
\mathrm{RunTail}_N(w,r) + (L2^{-m})^r + r\,b_{\delta,\xi}(N)
\right),
\end{equation}
and also
\begin{equation}
\label{eq:dense-inner-smallw-envelope}
p_w(\delta)
\;\le\;
\frac{w(w-1)}{N-1} + (L2^{-m})^w + w\,b_{\delta,\xi}(N).
\end{equation}
Consequently,
\begin{equation}
\label{eq:dense-inner-best-envelope}
p_w(\delta)
\;\le\;
\min\!\left\{
B_w^{\mathrm{run}}(\delta;\xi,m),\ 
(1-\varepsilon)^w + N2^{-m} + q_{\delta,\varepsilon}(N)
\right\}.
\end{equation}
\end{corollary}

\begin{proof}
The global bound \eqref{eq:dense-inner-global-envelope} is exactly the first part of
Theorem~\ref{thm:inner-dense-interface}.
The bound \eqref{eq:dense-inner-run-envelope} is the second part of
Theorem~\ref{thm:inner-dense-interface}, together with the exact run-tail identity
\eqref{eq:run-lower-tail-exact}.
The optimized form \eqref{eq:dense-inner-run-opt-envelope} is immediate by minimizing over $r$.
For the specialization \eqref{eq:dense-inner-smallw-envelope}, apply
\eqref{eq:dense-inner-run-envelope} with the choice $r=w$.
Then use \eqref{eq:runs-adjacent-bound} to bound
\[
\PP[R(U)<w]\le \frac{w(w-1)}{N-1},
\]
which yields \eqref{eq:dense-inner-smallw-envelope}.
The best-envelope form \eqref{eq:dense-inner-best-envelope} is then immediate.
\end{proof}

\paragraph{Why these are the right envelopes to optimize.}
The bound \eqref{eq:dense-inner-global-envelope} is designed for the linear-weight regime:
its first term decays exponentially in $w$, its second term exposes the ON$\to$OFF memory penalty $N2^{-m}$,
and its third term is a pure large-deviation term in $N$.
The run-aware bound \eqref{eq:dense-inner-run-opt-envelope} is the sharper tool for finite-length optimization:
it keeps the exact combinatorics of the interleaved run count,
it turns the OFF mechanism into the genuinely multiplicative term $(L2^{-m})^r$,
and it can be optimized over $r$ separately for each weight $w$.
The simpler bound \eqref{eq:dense-inner-smallw-envelope} is a convenient specialization for
the regime where random interleaving typically creates nearly $w$ runs.

\begin{theorem}[Dense-inner linear-weight exponent]
\label{prop:dense-inner-linear-exponent}
Fix constants $\delta\in(0,1/2)$, $\xi>0$, and $\theta\in(0,1)$.
Let
\[
L:=\lceil(2+\xi)\delta N\rceil,
\qquad
m=\gamma\log_2 N
\]
with $\gamma>1$, and define
\begin{equation}
\label{eq:dense-inner-binomial-exponent}
E_{\mathrm{bin}}(\delta,\xi)
\;:=\;
(2+\xi)\delta\left(1-H_2\!\Bigl(\frac{1}{2+\xi}\Bigr)\right).
\end{equation}
For a fixed weight fraction $\omega\in(0,1)$ satisfying $\theta<1-\omega$, let
\[
w=\lfloor \omega N\rfloor,
\qquad
r=\lfloor \theta w\rfloor.
\]
Then the dense inner satisfies
\begin{equation}
\label{eq:dense-inner-linear-exponent}
p_w(\delta)
\;\le\;
N^{O(1)}\,
2^{-N E_{\mathrm{in}}(\omega;\theta,\delta,\xi)},
\end{equation}
where
\begin{equation}
\label{eq:dense-inner-linear-rate-function}
E_{\mathrm{in}}(\omega;\theta,\delta,\xi)
\;:=\;
\min\!\left\{
\Psi_{\mathrm{run}}^{(w)}(\omega,\theta),\
E_{\mathrm{bin}}(\delta,\xi)
\right\},
\end{equation}
and $\Psi_{\mathrm{run}}^{(w)}(\omega,\theta)$ is defined in
\eqref{eq:run-rate-function-alpha}.
\end{theorem}

\begin{proof}
Apply \eqref{eq:dense-inner-run-envelope} with the choice $r=\lfloor \theta w\rfloor$.
The run-deficit term is controlled by Corollary~\ref{cor:run-lower-tail-alpha}:
\[
\mathrm{RunTail}_N(w,r)
\;\le\;
N^{O(1)}\,2^{-N\Psi_{\mathrm{run}}^{(w)}(\omega,\theta)}.
\]

Next,
\[
L2^{-m}
\;=\;
\Theta(N)\cdot N^{-\gamma}
\;=\;
N^{1-\gamma+o(1)},
\]
so since $r=\Theta(N)$ and $\gamma>1$,
\[
(L2^{-m})^r
\;=\;
2^{-\Omega(N\log N)}.
\]

For the binomial term, note that
\[
\frac{\delta N}{L}
\;=\;
\frac{1}{2+\xi}+O\!\left(\frac{1}{N}\right).
\]
Since $1/(2+\xi)<1/2$, a standard Stirling/Chernoff bound gives
\[
\PP[\Binomial(L,\tfrac12)\le \delta N]
\;\le\;
N^{O(1)}\,2^{-L\left(1-H_2(1/(2+\xi))\right)}
\;\le\;
N^{O(1)}\,2^{-N E_{\mathrm{bin}}(\delta,\xi)}.
\]
Multiplying by $r=O(N)$ only changes the polynomial prefactor.

Combining the three terms in \eqref{eq:dense-inner-run-envelope} yields
\eqref{eq:dense-inner-linear-exponent}.
\end{proof}




\begin{corollary}[Geometric dense-inner envelope above a logarithmic cutoff]
\label{cor:dense-inner-geometric-cutoff}
Consider the dense time-varying recursive inner encoder with memory
\[
m \;:=\; \gamma \log_2 n
\qquad\text{for a constant }\gamma>1.
\]
Fix $\delta\in(0,1/2)$ and choose any fixed $\varepsilon\in(0,1)$.
Let $p_w(\delta)$ denote the slice-to-tail probability from \eqref{eq:dense-inner-global-envelope}.
Assume the dense-inner tail bound (proved earlier in this section) gives
\begin{equation}
\label{eq:dense-three-term-recall}
p_w(\delta)
\;\le\;
(1-\varepsilon)^w \;+\; n\,2^{-m} \;+\; q_{\delta,\varepsilon}(n),
\end{equation}
where $q_{\delta,\varepsilon}(n)\le 2^{-c n}$ for some $c=c(\delta,\varepsilon)>0$.

Then there exist constants $\rho\in(0,1)$, $a\ge 0$, and $\kappa>0$
(depending only on $\delta,\varepsilon,\gamma$) such that
for every fixed cutoff $w_0(n)\ge \kappa\log_2 n$,
for all sufficiently large $n$, and for all $w\ge w_0(n)$,
\[
p_w(\delta)\;\le\; n^{a}\,\rho^{\,w}.
\]
\end{corollary}

\begin{proof}
Let $\rho:=1-\varepsilon\in(0,1)$.
We bound each term in \eqref{eq:dense-three-term-recall} by a common expression $n^a\rho^w$ for all
$w\ge w_0(n)\ge \kappa\log_2 n$.

\paragraph{Term 1.}
$(1-\varepsilon)^w=\rho^w\le n^a\rho^w$ for any $a\ge 0$.

\paragraph{Term 2.}
Since $m=\gamma\log_2 n$, we have $n\,2^{-m}=n\cdot n^{-\gamma}=n^{1-\gamma}$.
Then for all $w\ge w_0(n)\ge \kappa\log_2 n$,
\[
\rho^{\,w}\ \le\ \rho^{\,\kappa\log_2 n}
\ =\ n^{-\kappa\log_2(1/\rho)}.
\]
Choose any constant $a$ such that
\[
a - \kappa\log_2(1/\rho)\ \ge\ 1-\gamma.
\]
Then $n^{1-\gamma}\le n^{a}\rho^{\,w}$ for all $w\ge w_0(n)$.

\paragraph{Term 3.}
Since $q_{\delta,\varepsilon}(n)\le 2^{-cn}$ and $\rho^w\ge \rho^n$ is exponentially small in $n$ as well,
we have $q_{\delta,\varepsilon}(n)\le n^{a}\rho^{\,w}$ for all sufficiently large $n$ after increasing $a$ by an absolute constant if needed.

Combining the three bounds yields $p_w(\delta)\le n^a\rho^w$ for all $w\ge w_0(n)$.
\end{proof}

\begin{corollary}[Constant-factor geometric envelope]
\label{cor:dense-inner-constant-geometric}
Fix \(\delta\in(0,1/2)\) and choose \(\rho\in(0,1)\) with \(2\delta<\rho\).
Let \(m=\gamma\log_2 n\), and fix any logarithmic cutoff \(w_0(n)\ge \kappa\log_2 n\).
If
\[
\gamma \;>\; 1+\kappa\log_2\!\Bigl(\frac{1}{\rho}\Bigr),
\]
then for all sufficiently large \(n\) and all \(w\ge w_0(n)\),
\[
p_w(\delta)\le 3\rho^w.
\]
\end{corollary}

\begin{proof}
Apply Lemma~\ref{lem:dense-three-term-tail} with \(\varepsilon=1-\rho\).
The first term is exactly \(\rho^w\).
Since \(2\delta<\rho\), the binomial tail
\[
q_{\delta,1-\rho}(n)
=
\PP\!\left[\Binomial(\lfloor \rho n\rfloor,\tfrac12)\le \delta n\right]
\]
is exponentially small in \(n\), hence \(q_{\delta,1-\rho}(n)\le \rho^w\) for all sufficiently large \(n\) and all
\(w\ge \kappa\log_2 n\).
Also,
\[
n2^{-m}=n^{1-\gamma}
\le
n^{-\kappa\log_2(1/\rho)}
\le
\rho^w
\]
for all \(w\ge w_0(n)\), where the first inequality uses
\(\gamma>1+\kappa\log_2(1/\rho)\).
Summing the three terms proves the claim.
\end{proof}

\begin{remark}[The ON$\to$OFF term and why it is harmless in the framework]
\label{rem:dense-off-harmless}
The term $n2^{-m}$ in \eqref{eq:dense-three-term-recall} corresponds to the event that, while ON, the output stream
contains a length-$m$ all-zero window, causing an ON$\to$OFF transition.
For dense taps the expected ON lifetime under zero input is $\Theta(2^m)$, hence this term is unavoidable in a pointwise bound.
However, the serial concatenation framework never uses $p_w(\delta)$ for $w<d_{\min}(\Cone)$,
and for $w\ge d_{\min}(\Cone)=\Omega(\log n)$ the geometric factor $\rho^w$ yields an additional $n^{-\Omega(1)}$
that easily absorbs $n2^{-m}=n^{1-\gamma}$ as soon as $\gamma>1$.
Thus the dense inner conforms to the single inner interface without weakening the end-to-end guarantee.
\end{remark}




\paragraph{Work complexity (why this is the ``dense'' baseline).}
Each step computes $\langle k_t,s_t\rangle$ with a \emph{dense} $m$-bit mask, i.e.\ $\Theta(m)$ taps per symbol.
Thus total inner work is $\Theta(Nm)=\Theta(N\log N)$ for $m=\gamma\log N$.

\paragraph{Design notes (how the parameters are meant to be used later).}
\begin{itemize}
\item The term $(1-\varepsilon)^w$ is the clean exponential-in-$w$ contraction for intermediate weights.
\item The term $N2^{-m}$ is the rare ON$\to$OFF event. In tight finite-$N$ bounds, it is better to \emph{budget}
its cumulative contribution in the weight sum (rather than carry it pointwise for every $w$).
\item The binomial tail $q_{\delta,\varepsilon}(N)$ provides the exponentially small-in-$N$ large-deviation behavior
needed for the linear-weight regime when union bounding over exponentially many outer codewords.
\item Lemma~\ref{lem:dense-forced-termination} replaces the pointwise $N2^{-m}$ by $(L2^{-m})^r$ for small $w$,
using the fact that a random interleaver typically spreads $w$ ones into $\approx w$ runs when $w\ll \sqrt N$.
\end{itemize}

% ============================================================
\subsection{TODOs for this dense-inner module}
\label{subsec:inner-dense-todos}

\begin{enumerate}
\item \textbf{Sharper run-count tails for $R(U)<r$ beyond the $r=w$ case.}
We now have the exact finite-$N$ formula \eqref{eq:run-lower-tail-exact}. The next sharpening is to convert it into a
clean entropy-style exponent for linear weights, so the integration section can optimize over $r$ analytically rather
than numerically.

\item \textbf{Replace raw run count by an episode-aware small-weight law.}
The current finite-length evidence suggests that the main optimism in the old small-$w$ proof comes from treating
distinct runs as if they necessarily generated distinct ON-episodes. In reality several runs may merge into the same
episode when the encoder stays ON across the intervening zero gaps. A better theorem target is therefore not a direct
promotion of the old $(L2^{-m})^r$ contraction, but an exact late-placement base together with a controlled
multi-episode correction.

The current proof ladder is now:
\begin{enumerate}
\item isolate the dominant \(R(U)=h\) slice and use the exact late-placement base there;
\item use Corollaries~\ref{cor:isolated-order-late} and \ref{cor:isolated-early-count} to control exact early-start
statistics on that isolated slice;
\item treat \(R(U)<h\) as a lower-order adjacency remainder;
\item prove the remaining correction on the isolated slice through a small explicit statistic, rather than through the
full run count.
\end{enumerate}

Corollary~\ref{cor:forced-cluster-count} already shows that deterministic short-gap merging is negligible in the
large-$k$, \(h\approx 5\text{--}10\) regime currently driving the first-moment estimates. So the remaining correction
is more likely to be an \emph{episode-length} effect (multiple early episodes that still die quickly) than a large
hidden collapse of many runs into a few forced clusters.

The one-run corollary suggests the first concrete position-sensitive candidate:
assign each early isolated start at position \(y\) the suffix-length penalty
\[
(N-y+1)2^{-m} + \PP[\Binomial(N-y+1,\tfrac12)\le \delta N].
\]
Current overlap checks show that this already explains the isolated slice well at the tiniest weights, but still
undershoots in the \(h\approx 5\text{--}7\) window. So the missing theorem ingredient is now even clearer:
not the exact placement law, but the interaction term between multiple early episodes on top of it.

A particularly clean next target is to express that correction through the hypergeometric early-start count from
Corollary~\ref{cor:isolated-early-count}, or even just through the exact expected early-pair count from
Corollary~\ref{cor:isolated-early-pair-mean}, rather than through the full run count. Current overlap fits also
suggest that the coefficients in such an early-pair correction vary smoothly with the one-run base parameter
\(q_1\) and with \(1/m\), which makes this look much closer to a real analytic theorem target than a one-off fit.

\item \textbf{The first one-episode correction is too weak.}
Lemma~\ref{lem:dense-late-start-one-episode} and Corollary~\ref{cor:dense-late-start-slice} give a very clean
theorem-shaped correction to the exact late-placement base, but current numerics show that this first-episode-only
bound is still far too loose to explain the observed small-weight regime. So the next real theorem target must exploit
suppression from \emph{multiple} early episodes, not just the first one.

\item \textbf{Explicit constants in the binomial tail.}
If we want explicit finite-$N$ numerics, replace the entropy bound by an explicit Chernoff/Bennett form
with stated constants.

\item \textbf{Budgeting the $N2^{-m}$ term.}
In the global sum over $w$, it is typically better to choose a cutoff $W_2$ and show
$\sum_{w\le W_2} A_w^o \cdot N2^{-m}$ is below a chosen budget $\eta$; this removes a nuisance polynomial term.

\item \textbf{Unify notation with the sparse+refresh inner.}
Both the dense and sparse-refresh inners use the same ON/OFF narrative; we should unify the definitions
(episodes, termination windows, and ``fair coin while ON'') so the sparse-refresh section reuses this scaffolding.
\end{enumerate}
